\chapter*{How to use this book}
\addcontentsline{toc}{chapter}{How to use this book}
\stuck{InfoEdge Data Scientist interview experience (Ayush Agarwal)}{\ytid{nWyBQhGnHaM}}

\section*{What this book is, and how it fits with your OA book}

You already have the \emph{OA question bank} (675 short questions for the 40-question online test).
That book trains speed. This book trains \textbf{understanding and speaking}: Rounds 1 and 2 are
``drilling'' rounds in which an interviewer asks three or four questions per topic and keeps asking
``why?'' until you either explain the idea from first principles or run out of road. Round 3 is a
live Python notebook. Round 4 is an open case study. None of these can be passed by remembering
formulas; all of them can be passed by \emph{understanding}.

So every chapter has two halves.
\begin{enumerate}
\item \textbf{Teach.} The concept is built from nothing. No prerequisite is assumed: if a chapter
  needs a derivative, a dot product or a probability, it says in plain words what that is at the
  point it is needed. Every idea is followed by a \emph{small example} you can do with a pen in under
  a minute, and an \emph{in practice} box that places it in a Naukri, 99acres, Jeevansathi or Shiksha
  setting. Red boxes list the mistakes that interviewers fish for.
\item \textbf{Interview questions.} Each question is answered the way a strong candidate answers aloud:
  \begin{itemize}
  \item \textbf{\color{accent}In one line} -- the sentence you say first, so the interviewer knows you know.
  \item \textbf{\color{accent}The full answer} -- the logic, step by step, with the ``because'' behind every claim.
  \item \textbf{\color{accent}Worked example} -- numbers or a concrete scenario, so the idea is checkable.
  \item \textbf{\color{accent}If the interviewer pushes} -- the next two or three questions they will ask, already answered.
  \item \textbf{\color{accent}Trap to avoid} -- the wrong answer that sounds right.
  \end{itemize}
\end{enumerate}

\section*{The tags}
\begin{itemize}
\item \textcolor{pyqc}{\faStar\ \textbf{PYQ}} -- a question reported by InfoEdge candidates
  (online test, interview rounds, or the coding round). They come first in each chapter. InfoEdge does
  not publish papers, so the wording is a faithful reconstruction of what was reported.
\item \textcolor{deepc}{\faBrain\ \textbf{Think deeper}} -- a thought-provoking or ``unconventional''
  question of the kind the InfoEdge hire in the experience video called out: \emph{can an SVM do
  regression? what do you do with count data? is logistic regression a GLM? can a random forest
  extrapolate?} You asked for these to be at least 15\% of all questions. This edition contains
  \textbf{\total{qdeep} such questions out of \total{qtotal}} in total, and \textbf{\total{qpyq}} PYQs.
\end{itemize}

At the foot of \emph{every page} there is a ``Stuck on this page?'' link to one specific video for the
topic on that page. Where a direct video link could not be verified, the link is a YouTube search for the
exact video title, which lands on that video as the first result.

\section*{The InfoEdge selection process, from the Pre-Placement Talk}

\begin{center}
\small
\begin{tabularx}{\linewidth}{>{\bfseries}l X l}
\toprule
Stage & What is tested & Time\\
\midrule
Pre-screening & CV shortlisting + online test: about 40 MCQs in 40 minutes to 1 hour, from ML, DL,
programming and statistics. Programming MCQs are ``quite easy'' (for example converting a string to an
int and spotting the right output). & 40--60 min\\
Round 1 & Technical: Maths \& Statistics (33.33\%) and Classical ML \& Deep Learning (66.66\%). & 30 min\\
Round 2 & Same split, deeper. Reported as the round that ``tests your ability to explain the architecture
of ML and DL models'' on pen and paper. & 45 min\\
Round 3 & Programming. Part 1, data handling (Python, NumPy, Pandas): reading data, column naming,
statistical parameters, group-by, basic visualisation, basic text processing. Part 2, ML basics
(scikit-learn): train--test split, training a model, performance testing. & 30--45 min\\
Round 4 & Case study on a real-world data science problem: data challenges, modelling and design choices,
model assessment, thought process and team fit. ``There is generally no single best solution.'' & 45--60 min\\
HR & Company knowledge, why InfoEdge, work model. & 15--20 min\\
\bottomrule
\end{tabularx}
\end{center}

\subsection*{What the 2026--27 talk added}
\begin{itemize}
\item \textbf{Emphasis:} ``a lot more emphasis will be given to ML and DL''; GenAI is evolving, so
  \emph{basic} LLM knowledge is expected, but ``if you don't know topics like quantisation it will be
  fine''. The work is ``largely ML based, with some knowledge of DL and GenAI''.
\item \textbf{No DSA}, unless the discussion with a very strong candidate goes there.
\item The interviewers want ``a young data scientist who actually wants to understand business''.
\item \textbf{A spam example was used to explain what they look for:} do you look at recall and not just
  accuracy, can you build the right dataframe, do you frame it properly as a classification or a
  regression problem. Chapter 11 and Chapter 27 are built around exactly this.
\item \textbf{Products shown:} the Naukri job posting engine (100 crore+ applies per year), seeker
  acquisition and reactivation, apply generation, search and recommendation, the job--CV matching engine,
  Resdex search (60K recruiters, 16 lakh searches and 30 lakh candidates a day) organised as
  \emph{query intelligence $\to$ retrieval + ranking $\to$ personalisation}, PremiumX hybrid retrieval
  (dense + sparse, LLM ranking features, explainability), the SilverSkAI taxonomy (50 lakh raw skill strings
  normalised to 70K), AI REX agentic hiring (mandate, sourcing, screening, outreach, calibration agents),
  the Data Factory (400+ models; salary engine, company engine), Talent Pulse analytics, the content feed,
  Jeevansathi match recommendation (content, activity and collaborative models), and 99acres ML-based
  telecalling (user requirement, user--project matching, user ranking, then calls to high-intent leads).
  Every one of these appears as a case study in Part VI.
\item \textbf{Tools named:} classifiers, ensembles, kernel methods, dimensionality reduction, feature
  learning, Bayesian inference; ANNs, RNNs/LSTMs, Transformers, CNNs, label attention networks, Siamese
  networks; PyTorch, Keras/TensorFlow, Spark; Transformers, GPT, BERT, open LLMs, TTS/ASR, fine-tuning,
  agentic AI.
\end{itemize}

\subsection*{What a recent hire reported (IIT BHU, InfoEdge DS)}
\begin{itemize}
\item \textbf{Statistics:} p-value vs alpha. \textbf{Probability:} Naive Bayes built up from Bayes'
  theorem, with intuition rather than formulas.
\item \textbf{ML:} PCA via SVD and the scree plot; t-SNE and perplexity; Fisher LDA (push class means
  apart, shrink within-class spread); a little on decision trees.
\item \textbf{DL (the heaviest):} vanishing and exploding gradients; weight decay; how a network learns
  (backprop); preventing overfitting; Xavier vs He initialisation; time-series models (RNN, LSTM, GRU and
  their internals).
\item \textbf{Linear algebra:} rank, eigen-things, null space, and a Markov chain question.
\item \textbf{Round 3:} an MBA placement dataset. He spotted \emph{target leakage} first (salary is only
  present for placed students), checked missing values, checked outliers with \code{describe()} and
  \emph{did not} blindly remove genuine extreme marks, did value counts, a correlation matrix, logistic
  regression, a confusion matrix, precision, recall, F1 and accuracy.
\item \textbf{Round 4:} features predicting whether a 99acres user will buy (price-band concentration,
  location of recent searches); linear regression assumptions and what to do under heteroscedasticity
  (he proposed a GLM; the interviewer expected transformations; both were accepted); ``you are sales head
  of Naukri, give three ways to grow revenue, no clarifying questions''.
\item \textbf{Depth in basic ML:} count data needs Poisson regression (a GLM); SVMs can do regression (SVR).
\end{itemize}

\section*{How to answer aloud: the five-beat structure}
Interviewers grade \emph{structure} almost as much as content. For any ``explain X'' question:
\begin{enumerate}
\item \textbf{Define in one sentence} what X is and what problem it solves.
\item \textbf{Give the intuition} with a picture or an everyday analogy.
\item \textbf{Give the mechanism}: the formula or algorithm, and say what each symbol means.
\item \textbf{Give a tiny example}, ideally with numbers.
\item \textbf{Give the trade-off}: when it fails, and what you would use instead.
\end{enumerate}
If you are not sure, say what you \emph{do} know and reason aloud. ``I have not used it, but from the name
and from how X works I would expect\dots'' is a strong answer. Silence is the only bad answer.

\section*{A four-week plan}
\begin{center}\small
\begin{tabularx}{\linewidth}{l X}
\toprule
Week 1 & Part I (Chapters 1--5) and Chapters 6--9. Say every ``In one line'' aloud.\\
Week 2 & Chapters 10--17. Re-derive every worked example on paper without looking.\\
Week 3 & Part IV (deep learning) and Chapter 23. Draw LSTM, Transformer and backprop diagrams from memory.\\
Week 4 & Part V in a real notebook (type every snippet) and Part VI with a friend acting as interviewer.\\
\bottomrule
\end{tabularx}
\end{center}
